\documentclass[11pt]{article}
\usepackage{fontspec,amsmath,amssymb,amsthm,hyperref}
\usepackage[margin=1in]{geometry}
\setlength{\emergencystretch}{3em}
\title{Canonical finite-particle comparison from primitive latent data}
\author{LangevinLean mathematical exposition}
\date{}
\begin{document}\maketitle
\noindent Source module: \texttt{ProximalLangevin/MeanFieldCanonical.lean}. Compile this proof document with XeLaTeX or LuaLaTeX. The accompanying HTML pairs the mathematical argument with the exact Lean statements and records the scope of the formalization.


Given a finite-second-moment latent law \(\mu\), deterministic \(\theta_0\), globally Lipschitz \(\nabla V\), and \(T\ge0\), the initialized canonical product source supplies concrete independent Brownian inputs and actual integral solutions for all \(N>0\).


\section{Discharge every scalar Brownian law and full-path independence}


The initialized source already has the true Brownian coordinate laws and independent full paths. The injective particle-coordinate selection on its finite product preserves both properties. Thus the parameter coordinates and every particle’s latent coordinates form one independent standard Brownian family, including the first-copy overlap resolved by distinct coordinate tags.




\textbf{Correspondence with the displayed declarations.}


\begin{itemize}

\item \texttt{initialized\_particleBrownianFamily}: Actual interacting-particle Brownian coordinates on the iid initialized canonical source: every coordinate is standard Brownian and all full paths are independent, including common-parameter and first-particle coordinates.
\end{itemize}


\noindent Lean declarations: \texttt{initialized\_particleBrownianFamily}.

\section{Construct the reference once, then every particle process}


Use the actual initial pair \((\theta_0,X_0)\) on \(\mu\otimes Q_B\), whose square integrability and deterministic parameter projection were proved. Substitute this source and the actual Brownian coordinate process into the Brownian comparison theorem. It constructs continuous \(Y,\theta\), a finite \(C\ge0\), and for every \(N>0\) a continuous actual particle solution \(Z^N\) with exact iid initialization. The equations retain latent noise \(\sqrt2 B^{x,i}\) and parameter noise \(\sqrt{2/N}B^\theta\), and prove

\[\sup_{0\le t\le T}\big(\|\Theta_t^N-\theta_t\|_2+\|X_t^{i,N}-X_t^i\|_2\big)\le\frac{C_T}{\sqrt N}.\]

The coefficient is selected before the quantifier over \(N\). This controls the supremum of \(L^2\) norms, without claiming an expected path supremum.




The primitive inputs are a latent probability law with finite second moment, a deterministic parameter, a globally Lipschitz joint score and a finite horizon. The initialized product source supplies all Brownian coordinate laws, their whole-path independence, and initialization/noise independence. Substitute these into the already proved Brownian comparison theorem to construct one mean-field reference, then a particle solution for every positive \(N\). The output retains both integral equations and the deterministic parameter equation alongside the bound. Its uniformity is \(\sup_{t\le T}\|Z_t-R_t\|_{L^2}\); no exchange of supremum and expectation is made.




\textbf{Correspondence with the displayed declarations.}


\begin{itemize}

\item \texttt{exists\_canonical\_meanField\_particle\_comparison}: From a latent probability law with finite second moment, construct the mean-field reference and actual interacting-particle processes for every positive particle count. The displayed constant is independent of that count; its existence follows from proved finite-horizon moments of the reference.
\end{itemize}


\noindent Lean declarations: \texttt{exists\_canonical\_meanField\_particle\_comparison}.

\end{document}
